% Chapter 3: Methodology

\chapter{Methodology}

This chapter sets out the experimental methodology. Section~\ref{sec:research-design} covers the research design and ethics. Section~\ref{sec:phase1} describes how the Poisoned Evaluation Dataset was built. Section~\ref{sec:phase2} specifies the baseline meta-evaluation framework and gives formal definitions of the primary and secondary metrics. Section~\ref{sec:phase3} specifies the two pre-filtering modules and the ablation studies behind the design choices. Section~\ref{sec:phase4} defines the end-to-end evaluation protocol and additional analyses. Section~\ref{sec:implementation} records the implementation and statistical analysis framework.

\section{Research design}
\label{sec:research-design}

The study is empirical and experimental, organised into four sequential phases. Phase~1 (Section~\ref{sec:phase1}) builds the Poisoned Evaluation Dataset by injecting controlled adversarial perturbations into an established multi-hop QA benchmark. Phase~2 (Section~\ref{sec:phase2}) characterises three current LLM judges on the resulting dataset to quantify baseline failure rates. Phase~3 (Section~\ref{sec:phase3}) develops two pre-filtering modules motivated by the failure modes observed in Phase~2: a passage-level statistical classifier and a triplet-level LLM classifier. Phase~4 (Section~\ref{sec:phase4}) measures the end-to-end effect of each filter on judge reliability, runs the justification and multi-judge disagreement analyses, and conducts targeted ablations.

\begin{figure}[H]
\centering
\includegraphics[width=0.95\textwidth]{images/OverallExperimentalPipeline.png}
\caption{Overall experimental pipeline, showing the four phases from dataset construction through baseline judge scoring, pre-filter design, and end-to-end evaluation.}
\label{fig:methodology-overall-pipeline}
\end{figure}

The order of these phases is intentional; Phase~3 is informed by the empirical failure modes observed in Phase~2 rather than by prior assumptions about how judges fail.

\subsection{Ethical considerations}

The work involves synthesising adversarial content for experimental purposes inside a sandboxed evaluation environment. All adversarial passages are generated only for controlled offline analysis and are not deployed against any production system. The attack constructions are taken from published literature \citep{zou2025poisonedrag} and are presented to inform defensive research. The Poisoned Evaluation Dataset will be released with documentation and intended-use guidelines. No personal data is involved.

\section{Phase 1: Dataset construction}
\label{sec:phase1}

\subsection{Base benchmark}

We use HotpotQA in its distractor setting as the base benchmark. HotpotQA satisfies three properties relevant to this study. It requires multi-hop reasoning over multiple context passages, which is representative of the retrieval scenarios in which RAG systems operate and are most vulnerable. It provides ground-truth supporting passage indices and distractor passages per question, which lets us label poisoned passage indices exactly for supervised pre-filter training and for the True/Survived decomposition introduced in Section~\ref{sec:true-survived}. Finally, it is the base benchmark used in the original PoisonedRAG evaluation \citep{zou2025poisonedrag}, which permits direct comparison with prior work on RAG security.

We sample 100 HotpotQA hard questions whose answers require two supporting passages. These 100 base questions form the foundation of all 2{,}400 evaluation triplets in the final dataset.

\subsection{Triplet construction}

Each evaluation instance is a triplet $(q, C, a^{*})$ where $q$ is a question, $C = \{p_{1}, p_{2}, \ldots, p_{n}\}$ is the retrieved context (a set of passages), and $a^{*}$ is the ground-truth answer. The original HotpotQA distractor context contains $n = 10$ passages per question: the two supporting passages $p_{s_{1}}, p_{s_{2}} \in C$ and eight distractor passages.

For each base question and each combination of injection type and noise level, we build a poisoned triplet by replacing a fraction of $C$ with adversarial passages. Let $\rho \in \{0.2, 0.4, 0.6, 0.8\}$ denote the noise level, representing the fraction of original passages replaced. For each $(q, \mathrm{injection}, \rho)$ combination, we also keep a paired clean triplet (original $C$ unchanged) as a within-question control. The dataset therefore contains:
\begin{equation}
N = N_{q} \times N_{\mathrm{inj}} \times N_{\rho} \times 2 = 100 \times 3 \times 4 \times 2 = 2{,}400 \text{ triplets}
\end{equation}
This yields 1{,}200 poisoned and 1{,}200 clean triplets, with each poisoned triplet matched to a clean control at the same $(q, \mathrm{injection}, \rho)$ slot.

For each poisoned triplet, the indices of replaced passages are recorded in a \texttt{poisoned\_passage\_indices} field, denoted $\mathcal{P} \subseteq \{1, \ldots, n\}$. This explicit per-passage labelling enables supervised pre-filter training and residual-aware analysis. The dataset fields are summarised in \hyperref[app:schema]{Appendix D}.

\begin{figure}[H]
\centering
\includegraphics[width=0.95\textwidth]{images/PoisonedDatasetConstructionPipeline.png}
\caption{Poisoned dataset construction pipeline, from HotpotQA base questions to paired clean and poisoned triplets with explicit passage-level poison indices.}
\label{fig:methodology-dataset-construction}
\end{figure}

\subsection{Injection types}

We implement three injection types of increasing adversarial sophistication, each probing a different judge failure mode.

\textbf{Type 1, random noise.} Let $\mathcal{D}_{\setminus q}$ denote the union of all passages associated with base questions other than $q$ in the HotpotQA pool. A subset $\mathcal{P} \subseteq \{1, \ldots, n\}$ of context indices is selected uniformly at random. The candidate donor pool is filtered by cosine similarity so that the noise is genuinely off-topic: only passages with $\cos(\mathbf{e}_{p}, \mathbf{e}_{q}) < 0.3$ are kept as candidates, and each selected position is replaced by a passage drawn uniformly from this filtered pool. If the cosine filter produces an empty pool for a given question, the unfiltered donor pool is used as fallback so that a replacement is always available. The replacement passages are topically incoherent with $q$ by construction and can be detected by embedding-distance signals at low cost.

\textbf{Type 2, adversarial fact rewriting.} A subset $\mathcal{P}$ is selected as in Type~1. For each selected passage $p_{i}$, a generation model produces a rewritten variant $\tilde{p}_{i}$ via a single LLM call conditioned on $p_{i}$ alone. The prompt does not condition on $q$ or $a^{*}$, making the rewrite a passage-local perturbation rather than a question-targeted attack. The resulting passage is fluent and topically coherent with the original, but factually misaligned. The generation prompt is given in \hyperref[app:attack-prompts]{Appendix B}.

A post-hoc cosine-similarity guard rejects rewrites that drift below $\cos(\mathbf{e}_{p_{i}}, \mathbf{e}_{\tilde{p}_{i}}) < 0.3$, falling back to the original passage in those cases. A small fraction of generations that drift entirely off-topic are filtered out at construction time. Detecting Type~2 attacks requires factual reasoning rather than statistical anomaly detection.

\textbf{Type 3, PoisonedRAG-style injection.} Following the spirit of the black-box attack of \citet{zou2025poisonedrag}, each malicious passage is engineered to support a target wrong answer $a^{\dagger} \neq a^{*}$ while remaining topically coherent with the question. The construction is two-stage:

\textit{Stage 1, target wrong-answer generation.} A generation model is prompted with the question $q$ and the correct answer $a^{*}$ and asked to propose a plausible alternative that is specific, factually wrong, and could plausibly appear in source text.

\textit{Stage 2, malicious passage synthesis.} The generation model is given $q$, $a^{*}$, $a^{\dagger}$, and the original retrieved context, and is asked to produce a fluent passage that would mislead a downstream language model toward $a^{\dagger}$ rather than $a^{*}$. The exact prompt templates for both stages are given in \hyperref[app:attack-prompts]{Appendix B}.

\textit{Insertion mechanics.} For a triplet at noise level $\rho$, a number of malicious passages is determined and inserted into $C$ at uniformly-random positions, expanding the context length by that count. First, our malicious passage is the Stage~2 LLM output in isolation; we do not prepend $q$ as a separate retrieval sub-text, relying on topical coherence to minimize surface-level retrieval anomalies. Second, we insert the malicious passages \emph{additively} into $C$ rather than replacing existing passages. This preserves all original supporting and distractor passages, allowing us to isolate the effect of poison addition rather than poison-plus-replacement. The \texttt{poisoned\_passage\_indices} field records the positions of the inserted malicious passages in the resulting expanded context.

The additive insertion ensures that for every poisoned PoisonedRAG-style triplet, the original supporting passages are still present in the context, which is essential for the True/Survived decomposition introduced in Section~\ref{sec:true-survived}. 

\begin{figure}[H]
\centering
\includegraphics[width=0.95\textwidth]{images/PoisonedRAG-styleInjectionProcess.png}
\caption{PoisonedRAG-style injection process, showing target wrong-answer generation, malicious passage synthesis, and additive insertion into the retrieved context.}
\label{fig:methodology-poisonedrag-injection}
\end{figure}

\subsection{Dataset splits and statistics}

Triplets are partitioned into train, validation, and test splits with stratification by injection type and noise level to ensure balanced experimental conditions. The split sizes are detailed in Table~\ref{tab:splits}; the full dataset schema is provided in \hyperref[app:schema]{Appendix D}.

\begin{table}[h]
\centering
\caption{Dataset split sizes.}
\label{tab:splits}
\begin{tabular}{lrr}
\toprule
\textbf{Split} & \textbf{Size} & \textbf{Fraction} \\
\midrule
Train & 1{,}440 & 60\% \\
Validation & 480 & 20\% \\
Test & 480 & 20\% \\
\bottomrule
\end{tabular}
\end{table}

The pre-filter classifier is trained on the train split, hyperparameters are selected on validation, and all classifier-level performance reported is evaluated on the held-out test split.

Summary statistics for the full dataset:

\begin{table}[htbp]
\centering
\caption{Summary statistics of the Poisoned Evaluation Dataset.}
\label{tab:dataset-stats}
\small
\begin{tabular}{lp{0.55\linewidth}}
\toprule
\textbf{Statistic} & \textbf{Value} \\
\midrule
Base questions & 100 \\
Total triplets & 2{,}400 \\
Poisoned triplets & 1{,}200 \\
Clean triplets & 1{,}200 \\
Injection types & 3 (\texttt{random\_noise}, \texttt{adversarial\_fact}, \texttt{poisonedrag\_style}) \\
Noise levels & 4 (0.2, 0.4, 0.6, 0.8) \\
Total context passages scored & 25{,}587 \\
Maximum context length & 2{,}000 tokens \\
Maximum per-passage length & 500 tokens \\
Train / Val / Test split & 1{,}440 / 480 / 480 (60/20/20) \\
\bottomrule
\end{tabular}
\end{table}

\subsection{Quality assurance}

Each injection type passes type-specific quality checks before being added to the dataset.

For random noise, replacement passages are accepted only if their cosine similarity to the question is below $0.3$, so that the injected passages are genuinely off-topic.

For adversarial fact rewriting, each generated passage must satisfy three conditions: cosine similarity above $0.7$ with the original passage it replaces (to preserve retrievability), an automated consistency check confirming that the rewritten passage contradicts $a^{*}$, and an automated fluency check.

For PoisonedRAG-style, two operational conditions are verified: (i) the \emph{retrieval condition}, that the malicious passage appears in the top-$k$ retrieved results when injected into the knowledge base; and (ii) the \emph{generation condition}, that an LLM produces $a^{\dagger}$ when the malicious passage is provided as the sole context for the question.

A final integrity check confirms that no duplicate triplet IDs exist across splits and that the per-cell counts are balanced across $(\mathrm{judge}, \mathrm{injection}, \rho)$ combinations.

\section{Phase 2: Baseline meta-evaluation}
\label{sec:phase2}

\subsection{Judge selection}

We evaluate three current LLM judges spanning different provider and deployment profiles. 

\begin{table}[htbp]
\centering
\caption{LLM judges evaluated in this thesis.}
\label{tab:judges}
\small
\begin{tabular}{lll}
\toprule
\textbf{Judge} & \textbf{Model Identifier} & \textbf{Provider} \\
\midrule
GPT-5.4-nano & \texttt{gpt-5.4-nano}        & OpenAI           \\
Gemma        & \texttt{gemma-4-26b-a4b-it}  & Google AI Studio \\
DeepSeek     & \texttt{deepseek-chat} (V3.2) & DeepSeek        \\
\bottomrule
\end{tabular}
\end{table}

The models are chosen to cover different provider types rather than to exhaust the model space. Detailed configuration values are recorded in \hyperref[app:config]{Appendix E}.

\subsection{Scoring protocol}

All three judges receive an identical scoring prompt. The prompt provides the question $q$, the retrieved context $C$ (passages separated by blank lines, truncated at 2{,}000 tokens), and the answer $a^{*}$. The judge is asked to return three scores (faithfulness, answer relevance, and context relevance), each a real number in $[0, 1]$, in JSON format. The judge is not informed whether the context is poisoned, mirroring a zero-knowledge deployment scenario. The exact scoring and justification prompts are provided in \hyperref[app:evaluation-prompts]{Appendix A}.

We use the RAGAS conventions for the three metrics. \textbf{Faithfulness} measures whether the answer is supported by the provided context. \textbf{Answer relevance} measures whether the answer is on-topic for the question. \textbf{Context relevance} measures whether the retrieved context is relevant to the question. Faithfulness is the primary metric of interest because corpus poisoning compromises this dimension most directly.

\subsection{Primary metric: Faithfulness FPR}

Let $T_{p} = \{t_{1}, t_{2}, \ldots, t_{|T_{p}|}\}$ denote the set of originally-poisoned triplets and let $f(t)$ denote the faithfulness score assigned by a given judge to triplet $t$. We define a triplet-level \textbf{false positive} as a poisoned triplet that the judge scores at or above a threshold $\tau=0.5$:
\begin{equation}
\mathrm{FP}(t) = \mathbb{1}[t \in T_{p} \wedge f(t) \geq \tau]
\end{equation}

The \textbf{False Positive Rate (FPR)} of the judge is the fraction of poisoned triplets scored as faithful:
\begin{equation}
\mathrm{FPR} = \frac{1}{|T_{p}|} \sum_{t \in T_{p}} \mathbb{1}[f(t) \geq \tau]
\end{equation}

The threshold corresponds to the natural midpoint of the scoring interval and is the canonical operating point for binarising RAGAS scores. 

\subsection{Secondary metrics}

Three secondary metrics complement the headline FPR.

\textbf{Faithfulness gap} is the mean difference between clean and poisoned faithfulness scores:
\begin{equation}
\Delta_{f} = \overline{f}(T_{c}) - \overline{f}(T_{p})
\end{equation}
where $\overline{f}(\cdot)$ denotes the mean faithfulness score over the indicated triplet set and $T_{c}$ is the set of clean triplets.

\textbf{Stratified FPR} decomposes the headline FPR by injection type and noise level, which surfaces attack-class-specific failure patterns.

\textbf{Score distribution shape} characterises the judge's calibration regime. We report the fraction of scores in three bands: extreme low ($\leq 0.1$), middle ($0.1 < f < 0.9$), and extreme high ($\geq 0.9$). 

\subsection{Checkpoint and reliability}

All judge outputs are stored in a structured format that preserves triplet identifiers, judge identifiers, metric scores, poisoning labels, and filter condition labels. The output schema is summarised in \hyperref[app:schema]{Appendix D}.

\section{Phase 3: Pre-filter design}
\label{sec:phase3}

\subsection{Design criteria}

Both pre-filter modules are designed against four criteria: (i) pipeline integration as a drop-in component without modifying the judge prompt, (ii) computational tractability, (iii) multi-signal redundancy for the statistical filter, and (iv) configurable sensitivity.

\subsection{Statistical pre-filter: five signals}

The statistical pre-filter is a passage-level binary classifier. For each retrieved passage $p_{i} \in C$, the filter computes five signals, aggregated into a per-passage poison probability $\hat{y}_{i} \in [0, 1]$. Full configuration values are listed in \hyperref[app:config]{Appendix E}.

\textbf{Signal 1, embedding cosine anomaly.} This signal compares the passage embedding with the question embedding and flags unusually distant or unusually high-similarity passages.

\textbf{Signal 2, token-level entropy.} This signal uses language-model uncertainty to identify passages with unusual token-level probability patterns.

\textbf{Signal 3, fine-tuned classifier.} A lightweight transformer model is fine-tuned as a binary classifier on the training split to output a direct poison probability. 

\textbf{Signal 4, answer span recall.} This signal measures whether a passage preserves tokens associated with the ground-truth answer.

\textbf{Signal 5, cross-encoder relevance.} This signal scores query-passage relevance using an interaction-based cross-encoder.

\subsection{Signal aggregation}

The five signal scores form a feature vector $\mathbf{s}_{i} = (s_{1}, s_{2}, s_{3}, s_{4}, s_{5})$ per passage. We evaluate three aggregation methods on the validation split: weighted vote, XGBoost, and majority vote. 

The most effective aggregator on the validation set is selected as the primary aggregator for end-to-end experiments, with Weighted Vote retained as an unbiased baseline for ablation analysis. 

\begin{figure}[H]
\centering
\includegraphics[width=0.95\textwidth]{images/StatisticalPre-filterArchitecture.png}
\caption{Statistical pre-filter architecture, showing the five passage-level signals and their aggregation into a poison probability used for context filtering.}
\label{fig:methodology-statistical-filter}
\end{figure}

\subsection{Classifier backbone}

For Signal~3, we compare two encoder backbones: DeBERTa-v3-small and RoBERTa-base. Both are fine-tuned on the same train split with matched hyperparameters, listed in \hyperref[app:config]{Appendix E}.

\subsection{LLM pre-filter: triplet-level classification}

The second pre-filter strategy uses Mistral-small-latest as an LLM classifier. This model sits outside the families of both the generation models and the judges, reducing in-family recognition confounds.

Because HotpotQA is a multi-hop task, evaluating passages in isolation is not appropriate: legitimate supporting passages may only become useful in combination with other passages. We therefore adopt a triplet-level strategy. The LLM receives $q$, $a^{*}$, and the full context $C$ in a single prompt, and it flags the indices of any passages it identifies as adversarial. 

The output is a set of flagged indices $\hat{\mathcal{P}} \subseteq \{1, \ldots, n\}$. In deployment, a triplet for which $\hat{\mathcal{P}} \neq \emptyset$ is excluded from evaluation entirely, making this a \emph{content-preserving} filter, unlike the \emph{content-modifying} statistical filter. The two strategies operate at different abstraction levels and admit different failure modes, which is why we evaluate both.

\begin{figure}[H]
\centering
\includegraphics[width=0.95\textwidth]{images/LLMPre-filterStrategy.png}
\caption{LLM pre-filter strategy, showing triplet-level inspection of the question, answer, and full context before flagging suspected adversarial passages.}
\label{fig:methodology-llm-filter}
\end{figure}

\subsection{Filtered context construction}

For the statistical filter, the output is a modified context $C^{\prime} = C \setminus \{p_{i} : \hat{y}_{i} \geq \tau\}$ that the judge then scores.

For the LLM filter, the output is a triplet-level decision. We score the original context with the judge and compare verdicts on triplets the LLM filter flagged versus those it did not flag.

\section{Phase 4: End-to-end evaluation}
\label{sec:phase4}

\subsection{The True/Survived/Clean decomposition}
\label{sec:true-survived}

To evaluate the effect of a content-modifying filter, we must account for whether the original ground-truth supporting passages were overwritten by the attack. We decompose the originally-poisoned set $T_{p}$ using the explicit poisoning indices $\mathcal{P}$ and supporting passage indices $\mathcal{S}$. 

A triplet is in the \textbf{True} subset if the attack overwrote at least one supporting passage ($\mathcal{S} \cap \mathcal{P} \neq \emptyset$). It is in the \textbf{Survived} subset if the supporting passages remain intact. The two subsets partition $T_{p}$:
\begin{equation}
T_{p}^{\mathrm{true}} = \{t \in T_{p} : \mathcal{S}(t) \cap \mathcal{P}(t) \neq \emptyset\}, \qquad T_{p}^{\mathrm{surv}} = T_{p} \setminus T_{p}^{\mathrm{true}}
\end{equation}

Due to the additive nature of Type~3 (PoisonedRAG-style) injections, these attacks never overwrite original passages, meaning all Type~3 attacks inherently belong to the Survived subset by construction. Comparing baseline and post-filter performance within these respective subsets isolates the filter's effect on worst-case scenarios (True) versus scenarios where malicious content sits alongside original supporting passages (Survived).

\begin{figure}[H]
\centering
\includegraphics[width=0.95\textwidth]{images/ResidualAwareTrueSurvivedDecomposition.png}
\caption{Residual-aware True/Survived/Clean decomposition used in the end-to-end evaluation to separate support-overwritten cases from support-intact cases.}
\label{fig:methodology-true-survived}
\end{figure}

\subsection{Post-filter judge scoring}

Phase~4 re-runs each of the three judges on the filtered triplets under both filter strategies. The primary comparison metric is the change in FPR relative to baseline. 

\subsection{Justification analysis}
\label{sec:justification}

To probe whether judges can identify poisoning patterns when prompted appropriately, a small subset of triplets is re-scored using explicit poison-aware prompts. The poison-aware prompt is provided in \hyperref[app:evaluation-prompts]{Appendix A}.

\subsection{Multi-judge disagreement}

To test whether ensemble-based detection is a viable mitigation, we compute the standard deviation of faithfulness scores across the three judges for each baseline-condition triplet:
\begin{equation}
\sigma_{f}(t) = \mathrm{std}\bigl( f_{\mathrm{GPT}}(t), f_{\mathrm{Gemma}}(t), f_{\mathrm{DeepSeek}}(t) \bigr)
\end{equation}
We then compute the Pearson correlation between $\sigma_{f}(t)$ and the binary poisoning label $y(t) \in \{0, 1\}$. We additionally evaluate two ensemble decision rules at the operating threshold: a \emph{mean-score} ensemble and a \emph{majority-vote} ensemble.

\subsection{Statistical significance: McNemar test}

The end-to-end evaluation requires paired hypothesis testing because the same triplet is scored under multiple conditions. For each cell, we compute McNemar's test on the discordant pair counts:
\begin{equation}
\chi^{2} = \frac{(|n_{10} - n_{01}| - 1)^{2}}{n_{10} + n_{01}}
\end{equation}
The continuity correction is applied, and we treat $p < 0.05$ as significant.

\subsection{Ablation studies}

Two ablation studies are conducted on saved validation-split outputs. The \textbf{Per-signal ablation} removes each of the five signals in turn and re-evaluates the pipeline to observe individual signal contributions across injection types. The \textbf{Classifier backbone ablation} compares the two tested encoder backbones.

\section{Implementation and statistical analysis}
\label{sec:implementation}

The evaluation pipeline is implemented in Python using version-controlled configuration files. The implementation records raw judge outputs, pre-filter outputs, and analysis exports in structured files so that tables and figures can be regenerated from saved artefacts. The final configuration is summarised in \hyperref[app:config]{Appendix E}, and the reproducibility workflow is summarised in \hyperref[app:reproducibility]{Appendix G}.

Paired comparisons between conditions use McNemar's test for binary outcomes and the Wilcoxon signed-rank test for continuous scores. Unpaired comparisons across triplet subsets use the Mann-Whitney $U$ test. 

To account for inter-triplet dependencies stemming from the shared base questions, confidence intervals for classification metrics are computed via a question-clustered bootstrap resampling procedure (1{,}000 iterations). ROC curves are computed using scikit-learn's \texttt{roc\_curve} function, with the optimal operating point selected using the Youden index ($\mathrm{sensitivity} + \mathrm{specificity} - 1$).
